\input{preamble}

\begin{document}

\title{Research Journal}
\maketitle

\phantomsection
\label{section-phantom}

\tableofcontents

\section{Current objectives}
\label{section-objectives}

\begin{enumerate}
\item Keep working on Michele's deformation method in search
for a smoothing of $\text{SR}(\mathcal{M})$.

\item Try to understand the moduli space of generalised Kummer varieties
so that I get closer to the aim of the project as explained by Calla.
\end{enumerate}

\section{Stanley-Raisner schemes}
\label{section-sr}

\noindent
Here I follow \cite{luk1} for the basic definitions; these
are consistent with other sources.

Let $[n]:=\{0,\ldots,n\}$ be the set of all positive integers
from 0 to $n$, and let $\Delta_n$ denote the power set of $[n]$.
A subset $\mathcal{K} \subseteq \Delta_n$ is called a {\it simplicial complex} 
if for every $f \in \mathcal{K}$ all the subsets of $f$ are
also in $\mathcal{K}$. Elements of $\mathcal{K}$ are called {\it faces}.

Let $S=k[x_0,\ldots,x_n]$ be the polynomial ring in $n+1$ variables
over an algebraically closed field $k$.
If $a=\{i_1,\ldots,i_k\}\in \Delta_n$ we write $x_a \in S$ 
for the square-free monomial $x_{i_1}\cdot\ldots\cdot x_{i_k}$.
The {\it Stanley-Reisner} ideal
of a simplicial complex $\mathcal{K}$ is
$$
I_{\mathcal{K}}=\left<x_p : p \in \Delta_n\setminus \mathcal{K}\right>
\subseteq S.
$$
The {\it Stanley-Raisner ring} of $\mathcal{K}$ is
$A_{\mathcal{K}}=S/I_{\mathcal{K}}$
and the {\it Stanley-Raisner scheme} is $\text{Proj} A_{\mathcal{K}}$.

\section{Minimal resolutions}
\label{section-minimal-resolutions}

\noindent
In this section I explain what is a minimal resolution
and what are syzygies. The importance of this is that
our main tool to obtain flatness shall be
changesiArtin's criterion for flatness (\cite[Theorem 3.15]{Graffeo-Alessi}),
which requires computation of the first syzygy module.

\medskip\noindent
If an $S$-module $M$ is free, then it won't have any syzygies.
If $M$ is not free, then its first syzygy module
is the module generated by the relations which define  $M$.
More formally: consider a set $A$ of generators of $M$ 
and let $S^A$ be a free $S$-module generated by $A$.
This module maps onto $M$, and its kernel
is the {\it first syzygy module}:
$$
\xymatrix{
\text{Syz}^1(M)\ar[r]
&S^A\ar[r]
&M\ar[r]
&0.
}
$$
If $\text{Syz}^1(M)$ is free, it won't have
any syzygies. But if it is not free,
it will have syzygies of itself
(because by Hilbert Basis Theorem it is finitely generated),
giving
$$
\xymatrix{
\text{Syz}^2(M)\ar[r]
&S^{A_1}\ar[r]
&\text{Syz}(M)
&0
}
$$
where $S^{A_1}$ is the free $S$-module
generated by a set $A_1$ of generators of $\text{Syz}^1(M)$.

The {\it minimal resolution} of $M$ is
$$
\xymatrix{
&S^{A_k}\ar[r]
&\cdots\ar[r]
&S^{A_2}\ar[r]
&S^{A_1}\ar[r]
&M\ar[r]
&0.
}
$$
This means that the syzygy modules
{\bf do not} appear in the minimal resolution!
They are modules lying in-between each of the arrows
joining the free modules.
The minimal resolution is independent
(up to isomorphism) of the choices of generating
sets of each of the syzygy modules.

\subsection*{Explanation of twisting}
\label{subsection-explanation-twisting}
Consider
$$
\xymatrix{
0\ar[r]
&S(-1)\ar[r]^{\cdot x_0}
&S\ar[r]
&S/(x_0)\ar[r]
&0
}
$$
In $S(-1)$ constans have degree 1, linear polynomials have degree 2,
and so on.
This means that the map $\cdot x_0$ maps
degree-1 objects of $S(-1)$ to degree-1 objects in $S$;
i.e., it is degree-preserving.
Its image is $(x_0)$, so its cokernel is
$S/\text{img} = S/(x_0)$.

\subsection*{Explanation of M2 notation}
\label{subsection-m2-notation}
Usually we see in Macaulay something like
\[
\begin{array}{cccccc}
      & S      & \xleftarrow{} S^{16} & \xleftarrow{} S^{30} & 
	\xleftarrow{} S^{16} & \xleftarrow{} S \\
      &        &        &        &        &    \\
      & 0      & 1      & 2      & 3      & 4      \\
      & 1      & 16     & 30     & 16     & 1
\end{array}
\]
Here the module is again $S/I$,
but it doesn't appear in Macaulay's output.
Instead, $S/I$ is the cokernel of the first map $S^{16}\to S$,
and the map we had before $S \to M$ is the composition
$S^{16}\to S \to S/I$.

\section{Flat maps}
\label{section-flat-maps}

\noindent
In this section I start two examples of simple and intuitive flat maps:
deformations of parabolas and hyperbolas.

\medskip\noindent
To begin with I explain how flat morphisms in the category of schemes
are related to morphisms in the category of rings:
in practice we work at the level of rings 
and produce a flat map of rings $R \to S$
(namely, by Theorem \ref{theorem-grobner-flat-morphism}),
and then we just apply Spec to pass to geometry.

Typically  $R=k[t]$ is the parameter space
and $S=k[x_1,\ldots,x_n,t]$ is the affine (or projective)
coordinate ring with the extra variable $t$.
Then the fibres of $R \to S$ are the coordinate rings
of the fibres of the induced map $\Spec S \to \Spec R$.

In Sch,
$$
\xymatrix{
\Spec\kappa(\mathfrak{p})\times_{\Spec R}\Spec S\ar[r]\ar[d]
&\Spec S\ar[d]\\
\Spec\kappa(\mathfrak{p})\ar[r]
&\Spec R,
}
$$
while in Ring (applying the contravariant functor $\Spec^{-1}$
since the category of affine schemes is equivalent to Ring, 
Stacks Project \href{https://stacks.math.columbia.edu/tag/01I2}{01I2}
and \href{https://stacks.math.columbia.edu/tag/001J}{001J}),
$$
\xymatrix{
\kappa(\mathfrak{p})\otimes_RS&S\ar[l]\\
\kappa(\mathfrak{p})\ar[u]&R,\ar[l]\ar[u]
}
\qquad \text{that is,}\qquad 
\xymatrix{
R\ar[r]\ar[d]&S\ar[d]\\
\kappa(\mathfrak{p})\ar[r]&\kappa(\mathfrak{p})\otimes_R S.
}
$$
\medskip\noindent
Let's do two examples from \cite[II.3.3]{Hartshorne}.

\begin{example}
\label{example-flat-map-parabolas}
\begin{reference}
\cite[Example II.3.3.1]{Hartshorne}
\end{reference}
Let $k$ be an algebraically closed field, let
$$
X=\Spec k[x,y,t]/(ty-x^2),
$$
let $Y=\Spec k[t]$, and let $f:X \to Y$ be the morphism
determined by the natural homomorphism
$k[t] \to k[x,y,t]/(ty-x^2)$.
Then $X$ and $Y$ are integral schemes of finite type
(this means that the induced map of rings is of finite type)
over $k$, and $f$ is a surjective morphism.
We identify the closed points of $Y$ with elements of $k$.
For $a \in k$, $a \neq 0$,
the fibre $X_a$ is the plane curve 
$ay=x^2$ in $\mathbb{A}^2_k$, 
which is an irreducible, reduced curve.
But for $a=0$ the fiber is the nonreduced scheme given
by $x^2=0$.
Thus, the nonreduced scheme $x^2=0$ in $\mathbb{A}^2$ 
is a deformation of the irreducible parabola
$ay=x^2$ as $a \to 0$.

Algebraically, since $k[t]/(t) \cong k$,
the coordinate ring of the fibre is
$k[x,y,t]/(x^2-yt) \otimes_{k[t]}k \cong k[x,y]/(x^2)$.
The idea behind this isomorphism is that tensoring
with $k \cong k[t]/(t)$ is the same as ``killing $t$'',
namely, elements of the tensor product are polynomials
in $x,y,t$ but anytime $t$ appears it collapses to 0.
\end{example}

\begin{example}
\label{example-flat-map-hyperbolas}
\begin{reference}
\cite[Example II.3.3.2]{Hartshorne}
\end{reference}
Let $k$ be an algebraically closed field, let
$$
X=\Spec k[x,y,t]/(xy-t),
$$
let $Y=\Spec k[t]$, and let $f:X \to Y$ be the morphism
determined by the natural homomorphism
over $k$, and $f$ is a surjective morphism.
We identify the closed points of $Y$ with elements of $k$.
For $a \in k$, $a \neq 0$,
the fibre $X_a$ is the hyperbola
$xy=a$ in $\mathbb{A}^2_k$, 
which is an irreducible, reduced curve.
But for $a=0$ the fiber is the reducible scheme given
by $xy=0$.
Thus, the reduced scheme $xy=0$ in $\mathbb{A}^2$ 
is a deformation of the irreducible hyperbola
$xy=a$ as $a \to 0$.
\end{example}

In my search for an invariant that would tell me
whether GN and $\text{SR}(\mathcal{M})$ can be deformation-equivalent,
I gathered from 
\cite{Sernesi-Overview} the following important meaning for
``deformation theory'' (as opposed to studying flat families!):

\begin{quote}
``{\it Deformation theory} is the study of infinitesimal deformations
as a tool to understand the local structure of the moduli space.
The goal is to be able to describe the restriction
of the universal family to a small neighborhood of $m \in \mathcal{M}$,
or, more precisely, its restriction to the germ of $M$ at $m$.

What is interesting here is that we can study order and infinitesimal
deformations even though the functor $F$ is not representable or simply we don’t
yet know it is. This is the most frequent case. Such an investigation will
reveal the infinitesimal properties at $[X]$ of a yet unknown global structure
on $M$ which will be hopefully understood at a subsequent stage of the
investigation.  In order words it turns out to be possible and convenient to
separate the {\it global moduli problem} from the {\it local moduli problem},
and deformation theory studies the latter, with the purpose of constructing a
family of deformations of a given object parametrized by the spectrum of a local
ring, and having properties as close as possible to a universal property.''
\end{quote}

\medskip\noindent
A fun note: in \cite[Example 5.9]{jan2} it is shown that
the SR scheme of 600-cell is rigid! Also, the argument used
is not enough to confirm that $\text{SR}(\mathcal{M})$ is rigid.
And also, $\mathbb{P}(\mathcal{K})$ is never rigid.

\section{Deformation theory}
\label{section-deformations-}

\begin{definition}
\label{definition-deformation}
A {\it deformation} is… A {\it smoothing} is…
\end{definition}

\noindent
The key fact on why the tangent space of the Hilbert scheme
is the set of first order deformations is the following
(\cite[Exercise 3.8]{Graffeo-Alecci}).

\begin{exercise}
\label{exercise-tangent-space}
Let $X$ be a variety and let $p \in X$ be a closed point.
The tangent space $T_pX$ is isomorphic to the
following vector space
$$
T_pX \cong \{\varphi:(D_\varepsilon,(\varepsilon)) \to (X,p)\}.
$$
\end{exercise}

\noindent
Indeed, any tangent vector $(D_\varepsilon,(\varepsilon))\to (\text{Hilb},[X])$
corresponds to a flat family $Z \to D_\varepsilon$
by the fact the Hilbert scheme is representing
the functor of flat families!
And of course any flat family with base the dual numbers
is what we call an infinitesimal deformation.

\section{Deforming from the syzygies (aka Michele's method)}
\label{section-michele-method}

\noindent
In this section I explain Michele's method
for deformations based on Artin's flatness criterion 
(Theorem \ref{theorem-artin-criterion}).

The main idea is to add a deformation parameter
for every nonzero (degree-3, in this case) monomial
in the quotient $R/I$. 
The key fact about Artin's criterion is that
we must require that the syzygies of $I$ 
are also be valid in the deformed ideal.
So we compute the syzygies of $I$ and then
impose them to the deformed generators.
In practice, this means that we can
find some constraints on the deformation variables,
which actually helps with the computations.
The theory behind this is that we are actually
computing the module $\Hom(I,R/I)$.
See \cite[p. 346]{Sernesi} to find that in fact,
for a general embedding of schemes $X \subset Y$
the sheaf $\mathcal{I}/\mathcal{I}^2$ is defined
as the {\it conormal sheaf}, and its dual is
the {\it normal sheaf}:
$$
N_{X/Y}=Hom_{\mathcal{O}_X}(\mathcal{I}/\mathcal{I}^2,\mathcal{O}_X)
=Hom_{\mathcal{O}_Y}(\mathcal{I},\mathcal{O}_X).
$$
And we know
(I should be able to prove this by hand)
that first order deformations are parametrized
by sections of the normal bundle.
So, Michele's method is nothing else but a
computation of the normal sections,
equivalently, the first order deformations.

The next step is to apply Jacobian criterion to find smooth fibres.
Finally we need to check for obstructions,
since we know that not every first order deformation
is an actual deformation 
(but I do believe that the other way is true,
namely that every deformation comes
from a first order deformation).

\medskip\noindent
Consider the following simple example:
let $R=\mathbb{C}[x,y]$ and 
$I=(x,y)^2=(x^2,xy,y^2)$.
The syzygies in this case are
\begin{equation}
\label{equation-syzygies-example}
y(x^2)-x(xy)=0,\qquad y(xy)-x(y^2)=0,\qquad y^2(x^2)-x^2(y^2)=0,
\end{equation}
the last of which is trivial.

Then we define 
\begin{equation}
\label{equation-deformation-ring}
\mathbb{C}[x,y]\otimes\mathbb{C}[a_1,a_2,a_3,b_1,b_2,b_3,c_1,c_2,c_3],
\end{equation}
this will be the deformation ring.
Then we look for a family of the kind $Z \subset X \times B$
where $B$ is the base of the deformation;
such a family will be given by an ideal
$$
\tilde{I}=
(x^2+a_1+a_2x+a_3y,
xy+b_1+b_2x+b_3y,
y^2+c_1+c_2x+c_3y).
$$
Then we can compute the syzygies in $R/I$.
We get
\begin{equation}
\label{equation-syzygies-defomed-example}
\begin{aligned}
y(a_1+a_2x+a_3y)-x(b_1+b_2x+b_3y)&=0\\
a_1y-xb_1=0 & \iff a_1=b_1=0.
\end{aligned}
\end{equation}

Finally, the following theorem explains why
we worried so much about the syzygies
(\cite[Theorem 3.15]{Graffeo-Alecci}):

\begin{theorem}[Artin's criterion for flatness]
\label{theorem-artin-criterion}
Let $R=\mathbb{C}[x_1,…,x_n]$ be the polynomial ring
in $n$ variables and complex coefficients.
Let also $A$ be a local Artinian algebra of finite type over $\mathbb{C}$.
Let also $I \subset R$ and $I_A \subset R_A=R\otimes A$.
Suppose that there is an exact sequence
\begin{equation}
\label{equation-artin-criterion1}
\xymatrix{
R^\ell\ar[r]
&R^m\ar[r]
&R\ar[r]
&R/I\ar[r]
&0
}
\end{equation}
\noindent
and a complex
\begin{equation}
\label{equation-artin-criterion2}
\xymatrix{
R^\ell_A\ar[r]
&R^m_A\ar[r]
&R_A\ar[r]
&R_A/I_A\ar[r]
&0
}
\end{equation}

\noindent
such that the part
$$
\xymatrix{
R^m_A\ar[r]
&R_A\ar[r]
&R_A/I_A\ar[r]
&0
}
$$
is exact. Then, if
\ref{equation-artin-criterion2}$\otimes_A\mathbb{C}
=$\ref{equation-artin-criterion1},
the module $R_A/I_A$ is $A$-flat.
\end{theorem}

\noindent

Next:
``Do the Jacobian criterion on the whole ring
\ref{equation-deformation-ring}
and check if we get any subspace that depends
only on the parameters $a_i,b_i,c_i$.
This means that we have a space in the base.
If there is a vector space of this kind,
we can discard it.''

\medskip\noindent
Here's a quick overview of my latest experiment with $\text{SR}(o-o-o-o)$:

\begin{enumerate}
\item Syzygy analysis shows that the deformation polynomials
are
\begin{align*}
&a0*x0^2 + a1*x0*x1 + a2*x1^2 + x0*x2 + a3*x1*x2 + b5*x2^2 + b6*x2*x3\\
&c1*x0^2 + c2*x0*x1 + x0*x3 + b5*x2*x3 + b6*x3^2\\
&c1*x0*x1 + c2*x1^2 + c3*x1*x2 + c4*x2^2 + x1*x3 + c5*x2*x3 + c6*x3^2\\
\end{align*}

\item Jacobian criterion analysis shows that
\begin{align*}
&x0*x2+x1^2\\
&x0*x3\\
&x1*x3+x2^2
\end{align*}
\noindent
is a smooth fibre.

\item Artin criterion guarantees that
tensoring the ideal sheaf exact sequence of
$\text{SR}(o-o-o-o)$ with an Artinian algebra $A$
will give a flat map. But what is $A$?
According to the previous step, we know that
a good tangent direction for deformation is the line
$a_2=c_4$ in $\mathbb{A}^{12}$.
Such a line is given by the ideal
$I=(a_0,…,\widehat{a_2},…,a_6,b_1,…,b_6,c_1,…,\widehat{c_4},…,c_6,a_2-c_4)$.
When we intersect this line with a fat point at the origin,
given by the ideal $\mathfrak{m}^2=(a_0,…,a_6,b_0,…,b_6,c_0,…,c_6)^2$,
we obtain
\begin{align*}
\frac{\mathbb{C}[a_0,…,c_6]}{
I+\mathfrak{m}^2}&=
\frac{\mathbb{C}[a_2,c_6]}{
(a_2^2,c_6^2,a_2-c_6)}
=\frac{\mathbb{C}[a_2]}{(a_2^2)}
=\frac{\mathbb{C}[c_6]}{(c_6^2)}:=A.
\end{align*}

\noindent
Thus we may apply Artin criterion to get a flat $A$-module
$R_A/I_A=(\mathbb{C}[x_0,…,x_3]\otimes A)/(\text{polynomials in last step})$.

\item Since these varieties are curves,
$H^2(X,T_X)=0$, and thus any first-order deformation
is an honest deformation.
\end{enumerate}


\medskip\noindent
Here's a homeless note: The Reese algebra is finitely generated when the ideal 
is generated by squarefree monomials.



\section{Gröbner degenerations}
\label{section-grobner-degenerations}

\subsection*{Gröbner bases}
\label{subsection-Grobner-bases}

\noindent
Gröbner bases are a tool to make actual computations with varieties.
This subsection contains some basic concepts that culminate
in Theorem \ref{theorem-grobner-flat-morphism}
--- yes, it's probably better to skip all this and just understand
how the theorem might produce a smoothing.

\medskip\noindent
Here's a short summary on what are Gröbner bases and related concepts from
\cite{Sturmfelds-Groebner-bases}.

Let $k$ be any field and $k[\mathbf{x}]:=k[x_1,\ldots,x_n]$. The monomials in 
$k[\mathbf{x}]$ are denoted $\mathbf{x}^a=x_1^{a_1}\ldots x_n^{a_n}$ and
identified with lattice points $\mathbf{a}=(a_1,\ldots,a_n)\in \mathbb{N}^n$. A
total order $\prec$ on $\mathbb{N}^n$ is a {\it term order} if the zero vector
is the unique minimal element and $\mathbf{a}\prec \mathbf{b}$ implies
$\mathbf{a}+\mathbf{c}\prec\mathbf{b}+\mathbf{c}$ for all
$\mathbf{a},\mathbf{b},\mathbf{c}\in \mathbb{N}^n$.

Given a term order $\prec$, every nonzero polynomial $f\in k[\mathbf{x}]$ has a
unique {\it initial monomial}, denoted $\text{in}_{\prec}(f)$. If $I$ is an
ideal in $k[\mathbf{x}]$ then its {\it initial ideal} is the monomial ideal
$$
\text{in}_{\prec}(I):=\left<\text{in}_{\prec}(f):f \in I\right>
$$
A finite subset $\mathcal{G} \subset I$ is a {\it Gröbner basis} for $I$ with
respect to $\prec$ if $\text{in}_{\prec}(I)$ is generated by
$\{\text{in}_{\prec}(g):g \in \mathcal{G}\}$. It is called {\it reduced} if for
any two distinct elements $g,g'\in\mathcal{G}$, no term of $g'$ is divisible by
$\text{in}_{\prec}(g)$. 

The reduced Gröbner basis is unique for an ideal and a
term order, provided one requires the coefficient of $\text{in}_{\prec}(g)$ in
$g$ to be $1$ for each $g \in \mathcal{G}$ (why?). Starting with any set of
generators of $I$, the {\it Buchberg algorithm} computes the reduced Gröbner
basis $\mathcal{G}$.

\medskip\noindent
Fix a weight vector $\omega=(\omega_1,\ldots,\omega_n)\in\mathbb{R}^n$. For any
polynomial $f=\sum c_i \mathbf{x}^{\mathbf{a}_i}$ we define the {\it initial
form} $\text{in}_\omega(f)$ to be the sum of all terms
$c_i\mathbf{x}^{\mathbf{a}_i}$ such that the inner product $\omega\mathbf{a}_i$
is minimal. For an ideal $I$ we define the {\it initial ideal} to be the ideal
generated by all initial forms
$$
\text{in}_\omega(I):=\left<\text{in}_\omega(f):f\in I\right>
$$
This initial ideal need not be a monomial ideal (like the initial ideal for a
term order). But is will be whenever $\omega$ is chosen "sufciciently generic".
If in addition $\omega$ is non-negative, then, as we shall se,
$\text{in}_\omega(I)$ is an initial monomial ideal in the earlier sense.

Let $\omega\geq 0$ and let $\prec$ be an arbitrary term order. We define a new
term order $\prec_{\omega}$ as follows: for $\mathbf{a},\mathbf{b}\in
\mathbb{N}^n$, we set
$$
\mathbf{a}\prec_\omega\mathbf{b} :\iff
\omega\cdot\mathbf{a}<\omega\cdot\mathbf{b}
\text{ or }(\omega\cdot\mathbf{a}=\omega\cdot\mathbf{b}\text{ and }
\mathbf{a}\prec \mathbf{b}).
$$
\medskip\noindent

A {\it polyhedron} is a finite intersection of closed half-spaces in
$\mathbb{R}^n$. Thus a polyhedron $P$ can be written as
$P=\{\mathbf{x}\in\mathbb{R}^n:A\mathbf{x}\leq \mathbf{b}\}$ where $A$ is a 
matrix with $n$ columns and as many rows as there are hyperplanes to intersect 
(to visualize it think of the case of equality, which is an affine
plane, then the inequality gives a half-space. Do this for as many
hyperplanes/equations required). If $\mathbf{b}=0$, then there exist vectors
$\mathbf{u}_1,\ldots,\mathbf{u}_m\in \mathbb{R}^n$ such that
\begin{equation}
\label{equation-cone}
P=\text{pos}(\mathbf{u}_1,\ldots,\mathbf{u}_m):=
\{\lambda_1\mathbf{u}_1+\ldots+\lambda_m \mathbf{u}_m
:\lambda_i \in \mathbb{R}_+\}
\end{equation}
A polyhedron of as in Eq. \ref{equation-cone} is called a {\it cone}.
(I think this is the convex hull of $\mathbf{u}_i$ but extending by any positive
constant.)

The faces of a polyhedron in $\mathbb{R}^n$ can be defined as follows. Let
$\omega\in \mathbb{R}^n$, viewed as a linear functional. Define
 $$
\text{face}_{\omega}(P):=\{u \in P:\omega\cdot u\geq \omega v \;\forall v\in P\}
$$
Every subset $F$ of $P$ which has this form, that is, which maximizes some
linear functional $\omega$, is called a {\it face} of $P$.

If $P\subset \mathbb{R}^n$ is a polyhedron and $F$ a face of $P$, then the 
{\it normal cone}  of $P$ at $F$ is
$$
\mathcal{N}_P(F):=\{\omega \in \mathbb{R}^n:\text{face}_\omega(P)=F\}
$$
That is, the set of functionals $\omega$ for which $F$ is a face
$\text{face}_\omega(P)$.

A {\it polyhedral complex} $\Delta$ is a finite collection of polyhedra in
$\mathbb{R}^n$ such that
\begin{enumerate}
\item if $P \in \Delta$ and $F$ is a face of $P$, then $F \in \Delta$.
\item if $P_1,P_2\in\Delta$, then $P_1\cap P_2$ is a face of $P_1$ and of $P_2$.
\end{enumerate}
The support of a complex $\Delta$ is the union of all the elements of $\Delta$.
A complex $\Delta$ which consists of cones is called a {\it fan}. The collection
of normal cones $\mathcal{N}_P(F)$ as $F$ ranges over the faces of $P$ is a fan
called the {\it normal fan} of $P$.

The {\it Newton polytope} is the convex hull
of the indices $\alpha_i \in \mathbb{N}^n$ of the polynomial 
$f=\sum c_ix^{\alpha_i}$, where 
$x^{\alpha_i}:=x_1^{\alpha_i^1}\ldots x_n^{\alpha_i^n}$.

\begin{proposition}
\label{proposition-equivalence-classes-of-weight-vectors}
\begin{reference}
\cite[Proposition 2.3]{sturmfelds-groebner}
\end{reference}
Each equivalence class of weight vectors is a relatively open convex polyhedral
cone. This is implied by the fact that
$$
C[\omega]=\{\omega'\in \mathbb{R}^n:
\text{in}_{\omega'}(g)=\text{in}_\omega(g),\qquad \forall g \in \mathcal{G}\}
$$
for a reduced Gröbner basis $\mathcal{G}$ of $I$ with respect to $<_\omega$.
\end{proposition}

The union of these cones is a fan called the {\it Gröbner fan}.

The state polytope is a polytope whose vertices correspond to
initial ideals with respect to different monomial orderings/weight vectors.

\begin{theorem}
\label{theorem-groebner-fan-and-state-polytope}
\begin{reference}
\cite[Theorem 2.2]{sturmfelds-groebner}
\end{reference}
Let $I$ be a homogeneous ideal in $k[\mathbf{x}]$. There exists a polytope
$\text{State}(I) \subset \mathbb{R}^n$ whose normal fan
$\mathcal{N}(\text{State}(I))$ coincides with the Gröbner fan $\text{GF}(I)$.
\end{theorem}

\subsection*{Homogenization}
\label{subsection-homogenization}

Following \cite{lara},
the {\it homogenization of $f$} in $k[x_1,\ldots,x_n,t]$ is
$$
f^{h;\omega}
:=\sum c_ax^at^{\text{max}\{\omega\cdot b:c_b\neq 0\}-\omega\cdot a}
$$
And the {\it homogenization of $I$} is $(f^{h;\omega}:f \in I)$.

Notice that this is almost the same reasoning
we used in the previous section,
with the caveat that by the definition
of the homogenization of a polynomial,
the weight vector has different sign.

\subsection*{Gröbner degenerations and flat families}
\label{subsection-Grobner-degenerations-and-flat-families}

\noindent
Here is the main theorem used to construct a flat map from the
Gröbner base analysis.

\medskip\noindent
For an ideal $I \subset S=k[x_1,\ldots,x_r]$ de denote
its homogenization ideal with respect to a weight 
vector $\lambda$ (i.e. the ideal generated
by the homogeneization forms with respect to $\lambda$)
by $\tilde{I}$.
\cite[Theorem 15.17]{Eisenbud} is the following:

\begin{theorem}
\label{theorem-grobner-flat-morphism}
For any ideal $I \subset k[x_1,\ldots,x_r]$, the
$k[t]$-algebra $S[t]/\tilde{I}$ is free ---and thus flat---
as a $k[t]$ module. Futhermore, […].

Thus $S[t]/\tilde{I}$ is a flat family over $k[t]$ of quotients
of $S$ whose fiber over $0$ is $S/\text{in}_\lambda(I)$ and
whose fiber over any $(t-u)$ for $0 \neq u \in k$ is $S/I$.
\end{theorem}

That is, Gröbner bases give degenerations in the flat family sense.
This was found out via \cite{Lara-Grobner-degenerations}.


\subsection*{Twisted cubic example}
\label{subsection-twisted-cubic}

\noindent
Now I put everything in practice to construct a smoothing
of the SR scheme of the simplicial complex $\bullet-\bullet-\bullet-\bullet$
to the twisted cubic.

\medskip\noindent
I start by recalling basic properties of the twisted cubic.

\begin{definition}
\label{definition-twisted-cubic}
The {\it twisted cubic} is the image of the Veronesse map of degree 3,
\begin{align*}
\nu: \mathbb{P}^1 &\longrightarrow \mathbb{P}^3 \\
[u:v] &\longmapsto [u^3:u^2v:uv^2:v^3]=[x_0:x_1:x_2:x_3]
\end{align*}
\end{definition}

\begin{lemma}
\label{lemma-twisted-cubic-is-variety}
\begin{reference}
\cite[Proposition 6.1]{Eisenbud-syzygies}
\end{reference}
The twisted cubic is an algebraic variety whose ideal is generated by the minors
of the matrix
$$
M_{3,1}=\begin{pmatrix}
x_0 & x_1 & x_2\\
x_1 & x_2 & x_3
\end{pmatrix}
$$
\end{lemma}

\begin{proof}
See \cite[Proposition 6.1]{Eisenbud-syzygies}.
\end{proof}

Thus, the ideal of the twisted cubic is generated by
\begin{equation}
\label{equation-twisted-cubic}
x_1^2-x_0x_2,\quad x_1x_2-x_0x_3,\quad x_2^2-x_1x_3
\end{equation}

The minimal resolution of the twisted cubic is
\begin{equation}
\label{equation-resolution-of-twisted-cubic}
\xymatrix{
0\ar[r]&\mathcal{O}(-3)^{\oplus 2}\ar[r]^{d_2}&
\mathcal{O}(-2)^{\oplus 3}\ar[r]^{d_1}&
\mathcal{O}_{\mathbb{P}^3}\ar[r]&\mathcal{O}_X\ar[r]&0
}
\end{equation}
where
\begin{equation*}
\label{equation-matrices-twisted-cubic}
d_1=\begin{pmatrix}
x_1^2-x_0x_2&x_1x_2-x_0x_3&x_2^2-x_1x_3
\end{pmatrix},\qquad 
d_2=\begin{pmatrix}
-x_2&x_3\\
x_1&-x_2\\
-x_0&x_1
\end{pmatrix}.
\end{equation*}
It may be checked that $d_1d_2=0$.

\medskip\noindent
Now I describe the Stanley-Reisner scheme of four points and three lines.

The twisted cubic should deform to the Stanley-Reisner scheme of 
$\bullet-\bullet-\bullet-\bullet$, which we denote by $X$.
 The ideal of this scheme is 
$(x_0x_2,x_0x_3,x_1x_3)$ inside $S:=k[x_0,x_1,x_2,x_3]$.

Its minimal resolution is
\begin{equation}
\label{equation-resolution-of-three-lines}
\xymatrix{
0\ar[r]&\mathcal{O}(-3)^{\oplus 2}\ar[r]^{d_2}&
\mathcal{O}(-2)^{\oplus 3}\ar[r]^{d_1}&
\mathcal{O}_{\mathbb{P}^3}\ar[r]&\mathcal{O}_X\ar[r]&0
}
\end{equation}
where
\begin{equation*}
\label{equation-matrices-three-lines}
d_1=\begin{pmatrix}
x_0x_2&x_0x_3&x_1x_3
\end{pmatrix},\qquad 
d_2=\begin{pmatrix}
0&-x_3\\
-x_1&x_2\\
x_0&0
\end{pmatrix}.
\end{equation*}

\medskip\noindent
Now I'll do a deformation by hand of the twisted cubic's minimal resolution.
(This was done before I new anything about homogenization formalism;
below I will do it again using the language introduced
in Section \ref{subsection-homogenization}.)

To do the deformation of the twisted cubic we replace
$$
x_0\rightsquigarrow t^ax_0,x_1\rightsquigarrow t^bx_1,
x_2\rightsquigarrow t^cx_2,x_3\rightsquigarrow t^dx_3
$$
So that the equations of the twisted cubic \ref{equation-twisted-cubic} become
\begin{align*}
\label{equation-twisted-cubic2}
\begin{aligned}
t^{2b}x_1^2&-t^{a+c}x_0x_2\\
t^{b+c}x_1x_2&-t^{a+d}x_0x_3\\
t^{2c}x_2^2&-t^{b+d}x_1x_3
\end{aligned}
\end{align*}
To arrive at the equations of $X$ we must have
\begin{equation}
\label{equation-inequalities}
2b>a+c,\qquad b+c>a+d,\qquad 2c>b+d.
\end{equation}
A valid solution is $a=1,b=2,c=2,d=0$.

We obtain deformation matrices

\begin{equation}
\label{equation-deformation-matrices-twisted-cubic}
\begin{aligned}
d_1'&=\begin{pmatrix}
t^4x_1^2-t^3x_0x_2 &
t^4x_1x_2-tx_0x_3 &
t^4x_2^2-t^2x_1x_3
\end{pmatrix}\\
d_2'&=\begin{pmatrix}
-tx_2 & x_3\\
tx_1 & -t^2x_2\\
-x_0 & t^2x_1
\end{pmatrix}
\end{aligned}
\end{equation}

From \cite[p. 10]{HarrMorr}: ``The {\it twisted cubics} --- rational normal
curves in $\mathbb{P}^3$ that have Hilbert polynomial $P_X(m)=3m+1$''

\medskip\noindent
Now let's have a look at the Gröbner fan of the twisted cubic's ideal using
gfan.

As computed in \texttt{gfan}, the Gröbner fan of the
twisted cubic's ideal is composed of 8 cones in 
$\mathbb{R}^4$, with a lineality space of dimension 2.
This means we can plot the fan in a plane.
One of these cones contains the vector $(5,-5,-5,5)$,
for which the initial ideal is exactly
$(-x_0x_2,-x_0x_3,-x_1x_3)$.



\medskip\noindent
Now I will use the language introduced in
Section	\ref{subsection-homogenization} to find adequate weight vectors.

To see how this works consider the negative of the 
weight vector we found in Section \ref{subsection-deformation-by-hand}:
$\omega=(-1,-2,-2,0)$.
Let's compute the homogenization of the equations of the twisted cubic
\ref{equation-twisted-cubic}. For this just observe that the dot products with
the exponents of each monomial of these equations are:
\begin{align*}
\omega\cdot(0,2,0,0)&=-4,\qquad \omega\cdot(1,0,1,0)=-3\\
\omega\cdot(0,1,1,0)&=-4,\qquad \omega\cdot(1,0,0,1)=-1\\
\omega\cdot(0,0,2,0)&=-4,\qquad \omega\cdot(0,1,0,1)=-2
\end{align*}
The important part is that the right column,
which corresponds to the monomials that generate
the ideal of the Stanley-Reisner scheme,
have larger weight and will give exponent zero
on the variable $t$ by definition of homogenization form.
We obtain the homogenization forms
\begin{align*}
x_1^2t-x_0x_2,\qquad x_1x_2t^3-x_0x_3,\qquad x_2^2t^2-x_1x_3.
\end{align*}

Now consider the weight vector $\omega=(5,-5,-5,5)$
 found with \texttt{gfan}:
\begin{align*}
\omega\cdot(0,2,0,0)&=-10,\qquad \omega\cdot(1,0,1,0)=0\\
\omega\cdot(0,1,1,0)&=-10,\qquad \omega\cdot(1,0,0,1)=10\\
\omega\cdot(0,0,2,0)&=-5,\qquad \omega\cdot(0,1,0,1)=0
\end{align*}
As long as the numbers in the right column are 
greater than those on the left column,
we are left with the same initial forms.
In this case we get the homogenization forms
\begin{align*}
x_1^2t^{10}-x_0x_2,\qquad x_1x_2t^{20}-x_0x_3,\qquad x_2^2t^5-x_1x_3,
\end{align*}

\section{Stanley-Raisner scheme of Grünbaum sphere and a determinantal CY3}
\label{section-grunbaum-sphere}

A Gröbner basis for the Stanley-Raisner ideal
of Grünbaum-Shreedharan's strange complex $\mathcal{M}$ is
\begin{align*}
x_6x_7x_8, x_4x_6x_8, x_3x_7x_8, x_3x_5x_7, x_3x_4x_8, x_2x_7x_8, \\
x_2x_5x_7, x_2x_5x_6, x_2x_4x_7, x_2x_4x_6, x_1x_4x_6,\\
x_1x_4x_5, x_1x_3x_8, x_1x_3x_6, x_1x_3x_5, x_1x_2x_5
\end{align*}

\noindent
as computed in gfan and also GAP, and also as shown in \cite{luk1}.

There is a Calabi-Yau threefold in $\mathbb{P}^7$ of degree
20 that could be the smoothing of $\text{SR}(\mathcal{M})$.
(As far as I know there could exist others,
though they aren't very well-known.)
It was first studied by Gulliksen-Neg\aa ard \cite{Gulliksen-Negard}.
Among the papers I read often for this project, it is mentioned in
\cite[Table 1 and last paragraph]{CGKK-Aritmethically-Gorenstein-CY-in-P7} 
and  \cite{cascade}.
Additionally, this variety
has been studied in mirror-symmetry context, see \cite{Jockers}.

In trying to immitate the computations I did for the twisted cubic,
I would like to find the equations of the GN variety.
In \cite{cascade} we see 
its description as
the ``$3\times 3$ minors of $4\times 4$ matrix with linear forms''.

Back in May 19, 2025 I computed what might be ideal of the GN variety.
That is, as the $3 \times 3$ minors of the matrix
\begin{equation}
\label{equation-matrix}
\begin{bmatrix}
x_1 & x_2 & x_3 & x_4 \\
x_5 & x_6 & x_7 & x_8 \\
x_2 & x_3 & x_4 & x_5 \\
x_6 & x_7 & x_8 & x_1
\end{bmatrix}
\end{equation}
Those minors are
\begin{align*}
&x_3^2 x_5 - x_2 x_4 x_5 - x_2 x_3 x_6 + x_1 x_4 x_6 + x_2^2 x_7 - x_1 x_3 x_7, \\
&-x_3 x_6^2 + x_3 x_5 x_7 + x_2 x_6 x_7 - x_1 x_7^2 - x_2 x_5 x_8 + x_1 x_6 x_8, \\
&-x_3^2 x_6 + x_2 x_4 x_6 + x_2 x_3 x_7 - x_1 x_4 x_7 - x_2^2 x_8 + x_1 x_3 x_8, \\
&x_4 x_6^2 - x_4 x_5 x_7 - x_3 x_6 x_7 + x_2 x_7^2 + x_3 x_5 x_8 - x_2 x_6 x_8, \\
&x_3 x_4 x_5 - x_2 x_5^2 - x_2 x_4 x_6 + x_1 x_5 x_6 + x_2^2 x_8 - x_1 x_3 x_8, \\
&-x_1 x_2 x_5 + x_1^2 x_6 - x_4 x_6^2 + x_4 x_5 x_7 + x_2 x_6 x_8 - x_1 x_7 x_8, \\
&-x_1 x_2^2 + x_1^2 x_3 - x_3 x_4 x_6 + x_2 x_5 x_6 + x_2 x_4 x_7 - x_1 x_5 x_7, \\
&x_1 x_3 x_5 - x_1 x_2 x_6 + x_5 x_6^2 - x_5^2 x_7 - x_3 x_6 x_8 + x_2 x_7 x_8, \\
&x_4^2 x_5 - x_3 x_5^2 - x_2 x_4 x_7 + x_1 x_5 x_7 + x_2 x_3 x_8 - x_1 x_4 x_8, \\
&-x_1 x_3 x_5 + x_1^2 x_7 - x_4 x_6 x_7 + x_4 x_5 x_8 + x_3 x_6 x_8 - x_1 x_8^2, \\
&-x_1 x_2 x_3 + x_1^2 x_4 - x_4^2 x_6 + x_3 x_5 x_6 + x_2 x_4 x_8 - x_1 x_5 x_8, \\
&x_1 x_4 x_5 - x_1 x_2 x_7 + x_5 x_6 x_7 - x_5^2 x_8 - x_4 x_6 x_8 + x_2 x_8^2, \\
&x_4^2 x_6 - x_3 x_5 x_6 - x_3 x_4 x_7 + x_2 x_5 x_7 + x_3^2 x_8 - x_2 x_4 x_8, \\
&-x_1 x_3 x_6 + x_1 x_2 x_7 - x_4 x_7^2 + x_4 x_6 x_8 + x_3 x_7 x_8 - x_2 x_8^2, \\
&-x_1 x_3^2 + x_1 x_2 x_4 - x_4^2 x_7 + x_3 x_5 x_7 + x_3 x_4 x_8 - x_2 x_5 x_8, \\
&x_1 x_4 x_6 - x_1 x_3 x_7 + x_5 x_7^2 - x_5 x_6 x_8 - x_4 x_7 x_8 + x_3 x_8^2
\end{align*}
Now in theory I would just do all I did for the twisted cubic,
but now for GN, find a nice weight vector in for which the 
initial ideal is the ideal of $\text{SR}(\mathcal{M})$, and win.

But several problems arise:
\begin{enumerate}
\item First I have to double check that ideal for GN. I'm not
sure if those $3 \times 3$ minors of the $4 \times 4$ matrix I used are
good. {\bf Edit:} I think it's right. At least
by skimming through my references on GN, which are:
\cite{Jockers, Jockers2, Bertin, Gross-moduli-abelian}, I don't find
a matrix exactly but the claim is that a ``generic'' matrix
of $4 \times 4$ with linear forms in the entries should do the job.
(So that ChatGPT matrix above has indeed linear forms as entries.)
But most importantly \cite{Bertin} put a betti tally like
the ones one computes in Macaulay2 and indeed it's the same
as the one I get.

\item Most importantly, and this is what happened that other time
I tried to do this (when by the way I was being led by
ChatGPT to do exactly what I'm doing now, only now I
know a little more what's going on), is that the Gröbner fan
of GN is not computable. It's too much data and there's no output.
So it's not clear how to find the weight vector.
\end{enumerate}

One thing to try: obtain more freedom
on the $4\times 4$ matrix, see if you can change
that and still find the same good candidate.

\medskip\noindent
A breakthrough discovery:
the minimal resolution of GN and $\text{SR}(\mathcal{M})$ is not the same.

After this observation (and now knowing if
minimal resolutions of fibres of flat maps are isomorphic),
I'm not sure how to proceed. Possible ways are:

\begin{itemize}
\item Keep trying to prove/find out somehow 
(maybe asking: Eduardo Esteves thinks it's not true,
and recommended that I asked someone at URFJ)
whether of not the minimal resolution claim is true.
I may also look for other invariants on the fibers of a flat map,
see Stacks Project 
\href{https://stacks.math.columbia.edu/tag/02NM}{02NM}.

\item Keep trying to find the smoothing,
but I exhausted my computational attempts to find a perfect pairing
between the sets of Gröbner bases,
as well as my attempts to get around the size
of the Gröbner fan and just compute it via gfan.
\end{itemize}

\noindent
{\bf So I think I'll focus instead on trying to deform
$\mathbb{C}P^2_9$ using gfan technology.}
Also I could read about Kullikov models for degenerations
and understand a little better about the strategy of smoothing these
combinatorial complexes to construct the special fibres
of degenerations of generalised Kummer 4-folds.

\medskip\noindent
As a note, in \cite{luk2}, where the table 
where GN was found lies,
it is claimed that there is  ``evidence''
that this table is complete.
The table counts arithmetically Gorenstein CY3s.
Adding up, either I find a smoothing 
that is not the table,
or prove that it doesn't exist.


\section{6-vertex triangulation of $\mathbb{R}P^2$}
\label{section-6-vertex-triangulation-of-RP2}

\noindent
Here I read \cite{KuhnelBanchoff1983}.

\medskip\noindent
The icosahedron quotiented by antipodal
map gives a triangulation of 6 vertices
of $\mathbb{R}P^2$.
This can be embedded in $\mathbb{R}^5$ 
inside the 5-simplex $\Delta^5$,
which has 6 vertices.
This embedding has the properties
of symmetry (icosahedral$/-\text{Id}$),
duality (three vertices of $\Delta^5$ determine
a triangle of $\mathbb{R}P^2_6$ if and only if
the remaining do not),
tightness (too complicated) and
secant tangency (a chord joining
two points of $\mathbb{R}P^2$ is contained
in a boundary 4-simplex of $\Delta^5$.

Then consider Veronesse embedding
$[x_0:x_1:x_2]\mapsto [x_0^2:x_1^2:x_2^2:x_0x_1:x_0x_2:x_1x_2]$.
Like $\mathbb{R}P^2_6$ is contained in the boundary of $\Delta^5$,
this embedding is contained in $S^4$.
It also satisfies properties of symmetry (of $\mathbb{R}^3$),
duality, tightness and secant-tangency.

Each of the third and fourth properties
characterizes $\mathbb{R}P^2_v$ among
the smooth embeddings of $\mathbb{R}P^2$
in $\mathbb{R}^5$ not lying in an affine hyperplane.
``Pohl and Kuipar have in fact shown that
any {\it topological} embedding of the real
projective plane into $\mathbb{R}^5$ not lying in 
a hyperplane which has the tightness property
alone must either be $\mathbb{R}P^2_v$ or $\mathbb{R}P^2_v$.''

With respect to the complex projective plane,
the situation has been less satisfactory.

\section{9-vertex triangulation of $\mathbb{C}P^2$}
\label{section-9-vertex-triangulation-of-CP2}



\medskip\noindent
{\bf September 25.}
Summary of Calla's talk.
The objective is to construct a deformation 
of hyperkähler generalized Kummer fourfolds
starting from a special fiber.
The special fiber is the smoothing
of a Stanley-Raisner variety.

Proving this result would 
``imply that a general hyperkähler generalized kummer
fourfold with polarization of degree 2 and fixed divisibility
is contained in 36 quartics''.

\noindent
{\bf Questions.}
\begin{itemize}
\item How did they know that this funny simplicial
complex gives the special fiber of the degeneration
upon smoothing??
\item How does that result about being contained
in 36 quadric follows?
\item How come the special fiber of a degeneration
of fourfolds is a surface?
\end{itemize}

\medskip\noindent
Now let's go over the details of Calla's
construction. In \cite{KuhnelBanchoff1983}
it was shown that there is a 9-vertex
triangulation $\Sigma$ of $\mathbb{C}P^2$.
This triangulation has 36 4-simplices.
Notice that the definition of SR scheme for
Calla might be different from ours: the {\it
Stanely-Raisner variety} of $\Sigma$ is the
union of 36 copies of $\mathbb{P}^4$, one
for each 4-simplex of $\Sigma$. Its ideal is
generated by 36 quartics.

Part 1. We denote by $\Sigma$ the
9 vertex triangulation of  $\mathbb{P}^2$,
and $\text{SR}_{36}$ its corresponding
Stanley-Raisner scheme.

\begin{enumerate}

\item Blow up 9 coordinate points in $\mathbb{P}^8$ 
$(1:0:\ldots:0),\ldots,(0:\ldots:0:1)$.
(Recall that by the construction of the
Stanley-Raisner scheme associated to
a simplicial complex, in this case with 9
vertices, the vertices of the SR scheme
are precisely the coordinate points of $\mathbb{P}^8$.)
This introduces 9 copies of $\mathbb{P}^7$.
Proper transform of  $\text{SR}_{36}$
intersects $\mathbb{P}^7$ in a union
of  20 $\mathbb{P}^3$s, denoted $T_{20}$.
This does not depend on which copy of $\mathbb{P}^7$ 
we look at.

As a note, recall that the proper
transform of a subvariety $Y \subset X$ 
is defined as the {\it closure} of
$\pi^{-1}(Y\setminus Z)$ where
the blow up occurs at $Z$. In this case
$Z$ is the 9 vertices and $Y=\text{SR}_{36}$.
Calla claims that
that the closure of the proper
transform of $\text{SR}_{36}$ 
intersects each of the exceptional divisors
(one $\mathbb{P}^7$ for every vertex)
in a union of 20 $\mathbb{P}^3$ 
regardless of which exceptional divisor we
look at. We call that intersection $T_{20}$,
and it is inside $\mathbb{P}^7$.

\item Now we turn to modify $T_{20}$ 
to satisfy Friedman conditions.
Blow up 8 coordinate points in $\mathbb{P}^7$.
The proper transform of $T_{20}$ intersects
$\mathbb{P}^6$ in a union of 10 $\mathbb{P}^2$ s,
which we call $S_{10}$.
This {\bf does} depend on which copy of
$\mathbb{P}^7$ and $\mathbb{P}^6$ 
we are looking at!
$S_{10}$ is a Stanley-Reisner variety
and in fact is a triangulation of a sphere.
(A polyhedron!?)

\item Now we proceed to smooth $S_{10}$.
It has seven vertices, which we blow up,
introducing seven copies of $\mathbb{P}^5$.
Each of the $\mathbb{P}^2$ s becomes
a del Pezzo 6.
Looking at the picture we
see that there are vertices with
different valency. For each vertex
where $d$ $\mathbb{P}^2$s meet,
we glue in a del Pezzo surface of 
degree $d$. 
We get $\tilde{S}_{10}$, which is a 
simple normal crossing surface,
with trivial canonical bundle
and satisfies $d$-stability 
Friedman conditions.

Then there exists a Kulikov model
of type III degenerations of K3s
(of degree $10$ in $\mathbb{P}^7$)
with special fiber $\tilde{S}_{10}$.
(This is local smoothing over a
small disc, does not give a projective
family of K3s.)
\end{enumerate}

{\bf Key move.}
$S_{10}$ is a triangulation of a sphere.
Blow up the singularities and look
how does the proper transform of $S_{10}$ 
intersect the exceptional divisors.
To produce a smoothing,
we glue in del Pezzo surfaces
at the exceptional divisors.

Now we shall do the same for $T_{20}$.
(Which is still not the blow up
of $\text{SR}_{36}$, but the intersection
of the proper transform of $\text{SR}_{36}$ 
with any of the 9 exceptional divisors.)
We will glue in our projecive degeneration of K3s.
Notice that Friedman only gave us existence
of local model, so we must find equations.

For now let's just say that there's a bunch of procedures and
the resulting $\tilde{T}_{20}$ is a normal crossing and 
has trivial canonical bundle.
Then we can run smoothing procedure and find equations of family
of CY3s in $\mathbb{P}^7$ by noting that $T_{20} \subset 10 \mathbb{P}^4$s.
We should find generators of the ideal of
the variety $10\mathbb{P}^4$
among minors of $3\times 5$ matrix, similarly to surface case.

\medskip\noindent
{\bf October 19.}
It looks as if Calla was defining the SR variety of the triangulation
in a different way to my definition in Section \ref{section-sr}.
Today I checked that the ideal generated by 36 quartics
(one for every facet in the triangulation)
is not the same as the SR ideal.
How to tell which definition is ``correct''?
Correct for what? To answer I must get better understanding
of the deformation process from the Kummer perspective,
but it's not clear where this idea that the SR scheme
corresponding to the 9-vertex triangulation
would give the special fiber of a degeneration of
Kummer fourfolds!

I also noticed today that in \cite{KuhnelBanchoff1983}
there is a mention of $\mathcal{M}$!
I think $\mathcal{M}$ might be the vertex figure of $\mathbb{C}P^2_9$!


\section{Calla's talk}
\label{section-calla-talk}

\noindent 
In this section I write literally
Calla's talk.

\medskip\noindent
Aim: construct a type III degeneration of
varieties of generalised Kummer type 4-folds
with polarisation degree 2, by starting from
a given special fibre and constructing a
smoothing.

Corollary. Proving this result will imply that a 
general HK variety of this type with polarisation
degree 2 (\& fixed divisibility) is contained
in 36 quartics.
(In particular, it is not contained in a cubic.)

\medskip\noindent
Reminder of what we mean by type III degen:

Figure with three strings of text:
\begin{itemize}
\item "Moduli space of gen. Kum. type
4-folds w/ fixed polarization" enclosed in a circle
in the center,
\item "Boundary of BB compactification
is made up of 1 dim components 
intersecting in dim 0" with an arrow
pointing at to lines intersecting the circle
of the last item, and
\item "Type III degen has special fibre corresponding
to this point", which has an arrow pointing
at the point of intersection of the last item.
\end{itemize}

\medskip\noindent
Method: 
\begin{itemize}
\item We give a candidate $\text{SR}_{36}$ for
the special fibre of this degen.
\item Since $\text{SR}_{36}$ is badly singular,
we make some modifications,
in order to be able to smoothe it.

\end{itemize}
\section{Moduli space of generalized Kummer type fourfolds}
\label{section-moduli-kummer }

\noindent
According to Calla's talk,
the whole point of constructing the smoothing
of $\text{SR}(\mathcal{M})$ is that it
should help us to smooth $\mathbb{C}P^2_9$,
which in turn should lie in the boundary
of the moduli space of generalized Kummer type fourfolds
with polarization degree 2.

Here's a quote from Calla's talk:
``Aim: construct a type III
degeneration of varieties of generalized
Kummer type 4-folds with polarization degree 2,
by starting from a given special fibre
and constructing a smoothing.''

\medskip\noindent
Let's address the following question.
Even if there exists a smoothing of $\text{SR}(\mathbb{C}P^2_9)$,
why would it be a generalized Kummer variety?

First we go over \cite[Theorem 0.2.1]{Fausk}:

\begin{theorem}
\label{theorem-fausk-021}
A smoothing, if it exists, of the Stanley-Reisner scheme
of a triangulation of the 3-sphere yields a Calabi-Yau 3-fold.
\end{theorem}

\begin{proof}
Two key steps: show using \cite[Theorem 6.1]{Bayer-Eisenbud}
that the canonical bundle of SR(sphere) is trivial.
Then show using Hartshorne that triviality of
canonical bundle stays the same under deformations.
\end{proof}

\noindent
Sure, that says that if I smooth $\text{SR}(M)$ 
then I have a CY 3fold. That's only the vertex-figure
of $\mathbb{C}P^2_9$.
Still, this looks far from saying that a smoothing
{\it of $\mathbb{C}P^2_9$} is a generalized Kummer 4fold.


\section{Dual complex, Hironaka theorem and Kummer type manifolds}
\label{section-}

\noindent
These are notes from 
\href{http://verbit.ru/MATH/TALKS/Dual-complex-ddc_IMPA-2026.pdf}{
a seminar by Misha Verbitsky}
on his work with Nikon K.

\begin{definition}
\label{definition-snc}
Let $X \subset Y$ be a divisor in a smooth manifold.
We say that $X$ has normal crossings if
each of its components is smooth, and in a sufficiently
small neighbourhood of every point of $X$ the manifold
 $Y$ has coordinate system such that $X$ is identified with a union of
coordinate hyperplanes.
It has {\it simple normal crossings} if it has normal crossings and 
for each collection of irreducible components of the divisor,
their intersection is connected (possibly empty).
\end{definition}

\begin{theorem}[Hironaka resolution]
\label{theorem-hironaka}
Let $X$ be a complex variety.
Then there is a sequence of blow-ups with
smooth centers (=variety we are blowing up)
which results in a smooth resolution
$\pi:\tilde{X} \to X$ and the exceptional set of $\pi$ is SNC.
\end{theorem}

\begin{theorem}[Hironaka resolution of a pair]
\label{theorem-hironaka-pair}
Let $X$ be a complex variety and $Y \subset X$ a subvariety.
Then there is a sequence of blow-ups with smooth centers
which results in a smooth resultion
and the exceptional set of $\pi$ is a SNC divisor
and the preimage of $Y$ is also a SNC divisor.
\end{theorem}

\begin{definition}[Stepanov, 2004]
\label{definition-dual-complex}
Let $X \subset Y$ be a SNC divisor.
The {\it dual complex} of $X$ is a simplicial complex
whose vertices are the irreducible components $X_1,…X_n$,
and the $k$-simplices correspond to
intersections of $k+1$ of the irreducible components.
A {\it boundary} of a $k$-simplex is a $k-1$-simplex
obtained as the intersection of the divisors except one.
\end{definition}

\begin{theorem}[Stepanov]
\label{theorem-stepanov}
Let $x \in X$ be an isolated singularity
and $\tilde{X}$ its smooth resolution such that
the preimage of $x$ is a SNC divisor.
Then the homotopy type of its dual complex is independent from
the choice of a resolution.
\end{theorem}

\noindent
Stepanov's theorem was generalized by
Arapura-Bakhtary-Wlodarczyk.

\begin{theorem}[ABW]
\label{theorem-abw}
Let $Y \subset X$ be a subvariety,
and $\tilde{X} \xrightarrow{\pi}X$ the resolution of the pair
$Y \subset X$.
Let $D$ be the exceptional set of $\pi$ (which is a SNC)
and $\tilde{Y}$ the set of all components
$D' \subset D$ such that $\pi(D')\subset Y$.
Then the homotopy type of the dual complex of $\tilde{Y}$
is independent from the choice of resoution.
\end{theorem}

\begin{definition}
\label{definition-kummer-type}
We say that $Y \subset X$ is of {\it Kummer type} 
if $b_1(C(X,Y))=0$ where $C(X,Y)$ is the dual complex
of this pair, and $b_1$ the rational Betti number.
\end{definition}

\begin{theorem}
\label{theorem-kummer-type}
Let $X \subset M$ be a complex subvariety of Kummer type,
and $G$ a finite group holomorphically acting on $M$
and preserving $X$.
Then $X/G \subset M/G$ is also Kummer type.
\end{theorem}

\medskip\noindent
Now we discuss $dd^c$-lemma and the dual complex.

\begin{theorem}
\label{theorem-}
Let $X \subset M$ be compact complex ($M$ is not necessarily compact).
Assume that $X$ is Fujiki class C
(i.e. bimeromorphic to Kähler) (normal) of Kummer type,
that is, $b_1(S)=0$ where $S$ is the dual complex
of the pair $X \subset M$.
If $\eta$ is an exact $(1,1)$-form,
then there exists a function $f$ defined
in a neighbourhood of $X$ such that $\eta=dd^cf$.
\end{theorem}

\begin{remark}
\label{remark-normality}
Normality of $X$ may not be used.
Recall that normal= when all meromorphic functions are bounded.
\end{remark}

\begin{proof}
\noindent
Step 1. Consider the resolution $\tilde{X} \subset \tilde{M}$
With some extra blow ups we may assume that all the 
components of $\tilde{X}$ are Kähler.
Then we apply $dd^c$-lemma on any of the components.

\medskip\noindent
Step 2. $dd^cf=0$ means that it is a sum of holomorphic
and antiholomorphic,
which means that $\ker dd^c$ are
constant functions. There's also maximum princple
which implies that $f$ restricted to $\pi^{-1}(x)$ for any
$x$ in $X$ is constant.

\medskip\noindent
Step 3. Choose a function $f_i$ which satisfies
$dd^c f_i=\eta|_{D_i}$ on each of these components.
Then we prove that this satisfies the cocycle condition,
and we thus from $\eta$ we obtain a canonical element 
$\mu_\eta \in H^1(S,\mathbb{C})$.
[I think this cocycle condition is used
to glue the $f_i$ to a global function].

\medskip\noindent
Step 4. Extending $f$ to a neighbourhood…
\end{proof}

\noindent
Then there is a statement about the $dd^c$ lemma not
holding if $X$ is not of Kummer type: there is a form
that is not $dd^c$-exact:

\begin{theorem}
\label{theorem-ddc-counterexample}
Let $X \subset M$be a compact, positive-dimensional
subvariety of a manifold $M$, which is not of Kummer type.
Assume that there exists a non-contractible path in the
dual complex of the resolution of the pair, $S$,
which visits a component $D_0 \subset \tilde{X}$
only once, and the projection of $D_0$ to $M$
is positive-dimensional.
Assume $\dim M=2$.
Then there exists an exact $(1,1)$-form $\eta$
on a neighbourhood of $X$ such that $\eta \neq  dd^cf$ for any 
function $f$ on $X$.
\end{theorem}

\medskip\noindent
For the proof of ABW theorem we need the following results,
which were proved in 2002 in D. Abramovish, K. Karu, K. Matsuky,
J. Wlodarczyk, Torication and factorization of birational maps.

\begin{theorem}[Weak factorization theorem]
\label{theorem-weak-factorization-theorem}
Let $\varphi: X \dashrightarrow X'$ be a bimeromorphic map
of compact complex manifolds and $U \subset X$, $U' \subset X'$
the complements to the exceptional sets.
Then $\varphi$ can be decomposed 
into a sequence of blow-ups and blow-downs with
smooth centers.
\end{theorem}

\begin{remark}
\label{remark-meromorphic}
So what does it mean to be meromorphic? It means that there is
$$
\xymatrix{
&  Z \subset X \times X'\ar[dl]\ar[dr]\\
X & & X'
}
$$
where projections are proper, dominant and 1-to-1. [Maybe something
else I couldn't write down.
\end{remark}

\noindent
We need a version with boundary:

\begin{theorem}
\label{theorem-weak-factorization-for-pairs}
Let $\varphi:X \dashrightarrow X'$ be a bimeromorphic map
of compact complex manifolds, and $U \subset X$, $U' \subset X'$
be complements to the exceptional sets.
Consider SNC divisors $D \subset X$, $D' \subset X'$
and let $D_1,…, D_n$ and $D_1',…,D_n'$
be their irreducible components.
Assume that $\varphi$ takes $X\setminus D_i$
to $X' \setminus D_i'$
and this bimeromorphic map is proper
(that is, the closure of its grap in 
$X \setminus D_i \times X'\setminus D'$
projects properly to the components $X \setminus D_i$
and $X'\setminus D_i$).
Then $\varphi$ can be decomposed into a sequence
of blow-ups and blow-downs with smooth
centers, and the exceptional sets of blow-ups
have SNC with the proper preimages of $\cup  D_i$.
\end{theorem}

\noindent
Now we can prove Theorem \ref{theorem-abw},
which says that the homotopy type of the dual complex
is independent from the resolution.

\begin{proof}[Proof of Theorem \ref{theorem-abw}]
\noindent
Step 1. Take to resolutions of the pair.
Then they are bimeromorphic.
If we can show that the homotopy type of the dual
complex stays the same in each of the steps of the 
weak factorization theorem, then we are done.

\medskip\noindent
Step 2. We analyse the different possible cases of center in the 
first step. The first case is easy: the dual complex remains the same.
The second case I didn't understand.
The third case boils down to gluing a simplex to the 
original dual complex.
\end{proof}

\noindent
Then there is a $G$-equivariant version of this theorem
and of the weak factorization theorem.

\medskip\noindent
Finally we sketch Theorem \ref{theorem-kummer-type}.
Using $G$-invariant Hironaka, we
resolve, take quotient, take product, resolve, etc…
This would give a sequence of dual complexes…

\medskip\noindent
Take Hilbert scheme of torus, projected to symmetric
product of elliptic curve…
$$
\xymatrix{
(T^n)^{[n]}\ar[d]\\ \text{Sym}E.
}
$$


\section{Dual complex of a degeneration}
\label{section-dual-complex}

\noindent
Here I review \cite{KLSV} to understand
the relation between $\mathbb{C}P_9^2$ 
and the moduli of generalized Kummer varieties.

``It turns out that the right class of singularities
that still allow the definition of a meaningful dual complex is dlt.
In other words, the correct context for defining and intrincis
dual complex associated to a degeneration is that of minimal dlt
degeneration. The minimal dlt model $\mathcal{X}/\Delta$ 
is not unique, but changing the model has no effect on $|\Sigma|$.''

``Relevant for us is the fact that associated
to a $K$-trivial [I guees, trivial canonical bundle]
degenertion $\mathcal{X}/\Delta$ there is an intrinsic
(depending only on $\mathcal{X}^* /\Delta^*$) topological space
that we call (following [MN15]) {\it the essential skeleton} 
$\text{Sk}(X)$ associated to the degeneration.
For a minimal dlf degeneration of $K$-trivial varieties,
the essential skeleton $\text{Sk}(\mathcal{X})$ 
can be identified with the topological realization $|\Sigma|$
of the dual complex (cf Nicaise-Xu).''

``The purpose of this section is to make some remarks
on the structure of the essential skeleton $\text{Sk}(\mathcal{X})$ 
for a degeneration of hyper-Kähler manifolds''.

``[…] we note that the $\text{Sk}(\mathcal{X})$ for a MUM degeneration
of Calaby-Yau $n$-folds is predicted to be the sphere $S^n$.
This is true in dimension 2 by Kulikov's theorem,
and in dimensions 3 (unconditional)
and 4 (assuming additionally that degeneration is normal crossings
by Kollár-Xu.''

``6.2.6. To conclude, mirror symmetry
(via SYZ conjecture and Kontsevich-Soibelman)
predicts that the essential skeleton $\text{Sk}(\mathcal{X})$
for a MUM degeneration is $S^n$ and respectively
$\mathbb{C}P^n$ for Calabi-Yau's and respectively
hyper-Kähler's.
The following result is a weaker version of this statement,
saying that it holds in a cohomological sense.''

\begin{theorem}
\label{theorem-skeleton-cohomology}
\begin{reference}
\cite[Theorem 6.14]{KLSV}
\end{reference}
Let $\mathcal{X}/\Delta$ be a minimal dlt
degeneration of $K$-trivial varieties.
Assume that the general fiber $X_t$ 
is a simply connected $K$-trivial variety.

Assume that the degeneration is maximal unipotent.
Then
\begin{enumerate}
\item $H^* (\text{Sk}(\mathcal{X}),\mathbb{Q})
\cong \prod_iH^* (S^{k_i},\mathbb{Q})
\times_jH^* (\mathbb{C}P^{l_j},\mathbb{Q}))$,
where $k_i$ represent the dimensions
of the Calabi-Yau factos and $2l_j$ 
the dimensions of the hyper-Kähler factos
in the Beauville-Bogomolov decomposition of the general fiber $X_t$.

\item Conversely, the cohomology algebra
of the skeleton $\text{Sk}(\mathcal{X})$ 
determines the type of the Beauville-Bogomolov
decomposition of $X_t$.
\end{enumerate}
\end{theorem}



\section{Degenerations of Hilbert schemes (notes on Calla's thesis)}
\label{section-degenerations-hilbert}

\noindent
Notes on Calla's thesis on
degenerations of Hilbert schemes after
\href{https://www.youtube.com/watch?v=Lf4NYTvnO9I&t=969s}
{this video}.

The aim is to construct ``good'' degenerations
of Hilbert schemes of points over semistable
degenerations of surfaces. This was motivated by
the study of the Hilbert schemes of points
on ``maximally degenerate K3 surfaces'',
i.e. K3 surfaces that are on the boundary of
the compactified moduli space of K3s.

Basic setup: a flat projective family
of surfaces $X \to C \cong \mathbb{A}^1$ over
an algebraically closed field of characteristic 0.
Consider Étale local model given by
$\Spec k[x,y,z,t]/(xyz-t)$.
The general fibers are smooth, and
the special fibre $X_0$ over $0 \in C$ 
has simple normal crossing singularity.
We picture this as the three axes $x,y,z$ in $X_0$ 
intersecting at the origin.

Let $X^0= X\setminus X_0$ and $C^0=C\setminus \{0\}$.
Consider the Hilbert scheme of each fiber.
We have removed the limit over 0.
The question is: how can we compactify the Hilbert scheme
$\text{Hilb}^n(X^0/C^0)$ after we removed that fiber?
How to add a ``good limit''?

Starting from a ``big'' deformation space,
namely a fiber product $X \times_{\mathbb{A}^1}\mathbb{A}^n$,
Calla produces $X[n]$ after a sequence of blow-ups,
and later takes a quotient using a group action.
Thus produces a stack $\mathfrak{X}$
that contains the modifications
on the singular fiber
as well as the original fibers.

\begin{theorem}[Calla]
\label{theorem-calla}
The stack of Li-Wu stable length in 0-dimensional
subschemes in $\mathfrak{X}$ is Deligne-Mumford
and proper over $C$.
\end{theorem}

\noindent
It turns out this is a specific example
where things work out better than expected,
essentially by the way $X[n]$ was constructed.
The general case is not so ``nice''
(see [Maulik-Ranganathan]).

\medskip\noindent
{\bf Note.} At some point the main focus
of Calla's work turned to the Hilbert schemes.
I wonder if the generalised Kummer fourfolds studied in the
$\mathbb{C}P^2_9$ topic are actually examples
of these Hilbert schemes, or if it is a genuine
change of subject.


\section{Skeletal polytopes}
\label{section-polytopes}

\noindent
{\bf Problem:} determine the subgroup of (2,3,12) that
uniformizes the surface in which the facet of my master's thesis
chiral polyhedron is embedded.

\medskip\noindent
My strategy for finding this subgroup (so far) is just
understanding hyperbolic reflections groups in general.
Based on 
\href{https://youtu.be/9VAOF8MfL3M?si=7FVGX40Q_mpZVmen}
{Borcherd's video on Vinberg's paper},
I think that any hyperbolic reflection group might be isomorphic to 
a subgroup of $\text{GL}_2(\mathbb{Z})$, which acts on $\mathbb{H}^2$
by Mobius transformations.
Is this true? Can I write any examples of triangle groups?

\medskip\noindent
Here are the basic definitions taken from my thesis
(paper?! Still waiting for an answer…):

A \textit{skeletal polyhedron} in $\mathbb{E}^4$ consists of \textit{vertices}
(points in $\mathbb{E}^4$), \textit{edges} (segments between vertices) and
\textit{faces} (cycles on the graph determined by the vertices and edges) such
that:
    
\begin{enumerate}
\item every edge belongs to two faces, 

\item the
graph determined by the vertices and edges is connected, 

\item every compact
subset of $\mathbb{E}^4$ meets finitely many edges, and 

\item the \textit{vertex-figure}, defined as follows, is a connected graph. 
For any vertex
$v$, the vertices of the vertex-figure are the neighbours of $v$ and the edges
are segments joining any two neighbours that are both in some face.
\end{enumerate}

\noindent
A \textit{skeletal 4-polytope} in $\mathbb{E}^4$ consits of vertices, edges,
faces and \textit{cells} (skeletal polyhedra), such that
    
\begin{enumerate}
\item every face belongs to two cells,
\item the graph determined by the vertices and the edges is connected,
\item every compact subset of $\mathbb{E}^4$ meets finitely many edges, and
\item the vertex-figure at every vertex is a skeletal polyhedron.
\end{enumerate}

\noindent
Hereafter we omit the term ``skeletal'' when the context permits it. If
$\mathcal{P}$ is a 4-polytope, the set of vertices, edges, faces and cells is a
partial order with respect to inclusion. We say two such elements are
\textit{incident} if they are comparable. This poset is called the
\textit{abstract polytope associated to $\mathcal{P}$}. As we shall see,
4-polytopes defined as above are realizations of abstract polytopes in the sense
of \cite{ARP} and Section \ref{section-abstract-4-polytopes}.
	
A \textit{flag} is a 4-tuple of incident vertex, edge, face and cell. Two flags
are \textit{adjacent} when they differ by only one element. We call two flags
\textit{0-adjacent} if they differ by a vertex, \textit{1-adjacent} if they
differ by an edge, \textit{2-adjacent} if they differ by a face and
\textit{3-adjacent} if they differ by a cell.
	
A \textit{symmetry} of $\mathcal{P}$ is an isometry of $\mathbb{E}^4$ that
preserves it set-wise. The group of symmetries of $\mathcal{P}$ will be denoted
by $G(\mathcal{P})$. We call $\mathcal{P}$ \textit{regular} if $G(\mathcal{P})$
acts transitively on the set of flags, and \textit{chiral} if $G(\mathcal{P})$
induces two orbits on flags so that adjacent flags are on different orbits. In
fact, whether $\mathcal{P}$ is regular or chiral, its symmetry group acts
transitively on the sets of vertices, edges, faces and cells.

\medskip\noindent

My master's work was mostly based on \cite{petcox}. This paper contains a list
of so-called quiral polyhedra on $\mathbb{R}^4=\mathbb{E}^4$. My master's thesis
is basically using one of them to construct a 4-polytope 
(in the skeletal-polytope sense) 
that remains quiral; owing its quirality to
the quirality of the 3-cell.

The 4-polytope in question has the following combinatorial data:
120 vertices, 720 edges, 300 faces and 12 facets 
(=codimension 1 faces, in this case facets are the 3-cells).

The facet itself is a quiral skeletal polyhedron with
48 vertices, 72 edges and 12 faces.
At every vertex there are three edges,
so that we can think of this graph as embedded in a Riemann surface
which is a quotient of the hyperbolic plane
by a subgroup of the triangle group $(2,3,12)$.
Since the Euler characteristic is $48-72+12=-12$,
this is a surface of genus 7.


\medskip\noindent
Based on some data I computed back in the day, the Euler
characteristic of the 3-cell of Roli's cube (the first and probably simplest
example) is $16-24+6=2$, so $\mathbb{P}^1$.

\medskip\noindent
Here is a rather unsuccessful attempt to prove the
existence of the Riemann surface using \cite{Lando} from July 2025:

We may check that such a Riemann surface exists by
\cite[Construction 1.3.20]{Lando}. The two conditions that a graph must satisfy
are that the graph be connected, and 
\begin{equation}
\label{equation-sigma-alpha-phi}
\sigma\alpha\varphi=\text{id}
\end{equation}
where
$\sigma$ is the rotation about the vertex (sommet),
$\alpha$ is the rotation about the edge (arête),
and $\varphi$ is the rotation about the face. 
Note that these are not the three
rotations in a base flag: once we apply a symmetry, the next one should be
conjugated by the first.

When considering the conjugation scheme, the Eq \ref{equation-sigma-alpha-phi}
becomes $\alpha\sigma\varphi=\text{id}$. This is because we first apply
$\sigma$, but to apply $\alpha$ correctly we must conjugate by $\sigma$, that
is, multiply on the left of $\alpha$ by $\sigma^{-1}$ and on its right by
$\sigma$. Visual arguments that are sufficient for discussion right now show
that there is no change in face with these operations, so that no conjugation on
$\varphi$ is necessary.

In terms of $R_0,R_1$ and $R_2$, the generating reflections of the symmetry
group of a regular polyhedron, we have $\sigma_0=R_2R_1$, $\alpha_0=R_0R_2$ and
$\varphi_0=R_1R_0$. Here the $0$ denotes the fact that these are the vertex,
edge and face transpositions about the base flag. Notice that the transposition
about the edge is not simply $R_0$, it is a 180-degree rotation; the three
symmetries are orientation-preserving. Eq. \ref{equation-sigma-alpha-phi}
becomes $(R_0R_2)(R_2R_1)(R_1R_0)=\text{id}$.

This barely shows that indeed
there exists a Riemann surface underlying every regular convex polyhedron.

For our chiral 3-cell we obtain… I failed to prove Eq.
\ref{equation-sigma-alpha-phi} in this case. I made other attempts using
permutation groups in GAP, but still no success. Are there other ways of
constructing the Riemann surface? What kind of result do I expect from this
exercise?

{\bf July 25.} I conclude that it is worth the detour to correctly define a
constellation from the combinatorial data I have from my master's thesis as a
way to find my Riemann surface. Perhaps after this it wouldn't be too hard to
find the corresponding Belyi function and study it.

\medskip\noindent
The associated genus $7$ Riemann surface may admit
a Belyi function inducing a chiral dessin in the sense that it would not be
definable over the reals (as suggested by Sergey in the annotations of my
tesina). This means that the Riemann surface would not be biholomorphic to its
``conjugate''.

Based on some data I computed back in the day, the Euler
characteristic of the 3-cell of Roli's cube (the first and probably simplest
example) is $16-24+6=2$, so $\mathbb{P}^1$.

\medskip\noindent
I'm interested in knowing whether the Riemann surface is minimal
in $S^3$, just because this is a topic of interest for the group. It appears
this is possible: in \cite[Theorem 2.1]{brendle} we find Lawson's result that
there are minimal surfaces in $S^3$ for every genus.

But how are these surfaces embedded in $S^3$? What method we would use to know
confirm whether they are minimal or not? Is this surface or the
3-manifold obtained by gluing the facets really interesting?




\medskip\noindent
Here are some notes I took while reading Conway's book on quaternions.

\begin{lemma}
\label{lemma-generators-SO3}
$\text{SO}(3)$ is generated by simple rotations, i.e. a rotation that fixes a
vector.
\end{lemma}

\begin{lemma}
\label{lemma-double-cover-of-SO3}
Every element of $\text{SO}(3)$ has the form $x \mapsto q^{-1}xq$ for some
quaternion $q$, and is a simple rotation. The map that takes $q$ to the map
$[q]:x \mapsto q^{-1}xq$ is a 2-to-1 homomorphism from the group of unit
quaternions to $\text{SO}(3)$.
\end{lemma}

\noindent
Compare this with Graham's talk:
unit quaternions are $S^3$, and they act on the imaginary quaternions
(which are $\mathbb{R}^3$) by
$$
\rho(\alpha)\beta:=\alpha \beta \overline{\alpha}.
$$
That is, we have a homomorphism $\rho:S^3 \to \text{SO}_r(\mathbb{R})$ 
with kernel $\{\pm \text{Id}\}$.
Geometrically, this means that unit quaternions are represented
as rotations in $\mathbb{R}^3$ whose rotation axis is
(probably) $\alpha$.

Apparently the set of unit quaternions is called {\it spin group}
$\text{Spin}(3)$, and this is surely just $\text{SU}(2)$.

The fact that the covering is 2-to-1 is because opposite quaternions yield the
same rotation.

\begin{definition}
\label{definition-binary-icosahedral-group}
The {\it binary icosahedral group} is the preimage of the icosahedral group
$I\cong A_5$ by
the covering described above.
\end{definition}

As an aside, the manifold $\text{SO}(3)/I$ is the infamous Poincaré homology
sphere, also known as Mukai-Umemura threefold in Algebraic Geometry (double
check this).

\section{Cayleypy}
\label{section-cayleypy}

\noindent
Watching \href{https://www.youtube.com/watch?v=RkmBwlSyhfA}{Sasha's talk}
to understand how Cayleypy works. I wonder if I could use it
for the skeletal polytope problem.

The heuristics are (I've copied this literally from the key slide):
\begin{enumerate}
\item (Random walks to create ``training set''.)
\label{item-random-walks}
Run $N$ random walks on a Cayley graph starting
from the identity and each making $K$ steps.
Store the obtained group elements (permutations)
and numbers of steps they were achieved,
i.e. create dataset of pairs $(p,k)$
where $p$ are permutations 
and  $k$ - integer number of steps it was achieved by random waslk.
Create millions-billions of such pairs - that is a train set for neural net.

\item (Train machine learning model.)
\label{item-training}
Training neural network to predict  ``$k$ from $p$'',
i.e. input for neural net - vector describing permutation;
output - single number $k$ which represents how many
steps we need to get there by random walk
(``diffusion distance'').

\item (Graph search procedure.)
\label{item-graph-search}
Starting from given element ``$g$'',
for each of its neighbours compute prediction of neural network - 
choose neighbours with the smalls predictions:
number of random walks, graph structure, neural net, training, and so on.
Cayley graphs for permutation and matrix groups - 
natural examples where Python code can have unified structure.
\end{enumerate}

\noindent
An important remark is that after Step
\label{item-training} the net can only make
{\it imprecise} estimations.
That's why we need Step \ref{item-graph-search},
where morally we: stand at a point, ask the model
what node it thinks is best, move there, ask again, and so on.
Also, this ``heuristic explanation''
is not what they really do - it's just the idea.

\medskip\noindent
Why do we need neural networks?
We cannot handle the data of all the positions:
it's just too much data ($10^{75} \sim $ number of atoms in the universe!).
So we give machine learning a smaller data set
and hope it will somehow extrapolate.
There's no theorem for this.

\medskip\noindent
{\bf November 2.}
Two interesting topics found in \cite{Cayleypy} are:
\begin{enumerate}
\item ``Transforming men into mice'' idea 
(that's actually the name of a paper, but for now just a recalling device)
that we can compute genomic distance from one chromosome to another.

\item Sergey's ``wild hypothesis if the quasi-polynomial lengths/diameters could
be Erhart quasi-polynomials of polygons for some hidden reason''.  This
addresses the following situation, as I explained to Vinicio, that ``En
particular (en mis palabras): 
\href{https://www.cube20.org/}{se demostró} que el diámetro (mayor cantidad de
movimientos necesarios para llegar al estado “resuelto”) del grupo del cubo de
Rubik es 20. El cubo de Rubik es un subgrupo del grupo simétrico. 
Problema abierto: encontrar el diámetro del grupo simétrico completo.''
But after a glance at the conjecture table in \cite{Cayleypy}
we see that in fact this problem depends on the choice of generators
for $S_n$.
\end{enumerate}

\noindent
More precisely, the biological part discusses at least
two operations on a DNA chain: reversal and transposition.
The first one means we flip the order of a given sequence
and the second one means we move a sequence from one place
of the whole string of DNA to another.
I think the task here is to write this information
in terms of a Cayley graph and study its group properties
for very large values of sequence length $n$.

Next time: start at 5.2 ``Growth of Reversals and Transposons
on binary cosets'' to understand what Cayleypy team has done on this
problem.

\section{Vertex algebras}
\label{section-vertex-algebras}

\noindent
Trying to figure out if the Charged fermions 
(=Dirac sea, infinite-dimensional Clifford algebra)
can be represented in the (co-)homology group
of some moduli space of sheaves since December 2025.

In fact, I think this might already have been ``done''
in Nakajima's work. Recall \cite[Theorem 8.13]{Nakajima-lectures}:

\begin{theorem}[Nakajima, Grojnowski]
\label{theorem-main}
The following relations hold,
Then
\begin{equation}
\label{equation-main}
\begin{aligned}
[P_\alpha[i],P_\beta[j]]
&=(-1)^{i-1}i \delta_{i,-j}
\left<\alpha,\beta\right>\text{Id}
\qquad \text{if }(-1)^{\text{deg}\alpha \cdot \text{deg}\beta}=1,\\
\{P_\alpha[i],P_\beta[j]\}
&=(-1)^{i-1}i \delta_{i+j,0}\left<\alpha,\beta\right>\text{id}
\qquad \text{otherwise,} 
\end{aligned}
\end{equation}
\noindent
where $\left<\alpha,\beta\right>=p_*(\alpha \cap \beta)$
for $p:X \to \text{pt}$ the unique morphism
from $X$ to the point pt.
\end{theorem}

Next step: understand NS using huybrechts.


\subsection{Clifford algebras}
\label{subsection-clifford-algebras}

\noindent
Clifford algebras appear to be the finite-dimensional
version of the Charged Fermion algebra that we
shall try to find represented in the cohomology of
some sheaf moduli space.

Here are a few basic facts about the finite-dimensional
Clifford algebras from \cite{Frenkel} and Misha's 2026
summer course on Seiberg-Witten invariants:

\begin{enumerate}
\item (Definition.) The {\it Clifford algebra} of a vector space $(V,q)$ 
equipped with a scalar product $q$
is the quotient of the tensor algebra
by the relation $x \otimes y+y \otimes x+q(x,y)1$.

In the language of  \cite{Frenkel},
it's just the algebra generated by $e_i$ 
(which can be thought of as the generators of $V$)
subject to the relation $e_ie_j+e_je_i=-2\delta_{ij}$ 
(so $q=\delta$).

\item ($\mathbb{Z}_2$-grading.) The Clifford algebra
is $\mathbb{Z}_2$-graded by ``number of factors in the product of any element''.

\item (Relation to exterior power.) The Clifford algebra is filtered
and the associated graded is $\Lambda^\bullet V$.
Also, the infinite-dimensional version of the Clifford algebra
has a representation (the Fock module, i.e. highest-weight?)
on some infinite-dimensional vector space (see \cite{Nakajima-lectures}).

\item (Jethro.) We can also construct the infinite-dimensional version
(called Charged-free fermions) by starting from some abstract stuff
$\varphi\mathbb{C} + \varphi^*\mathbb{C}$ and
then it's possible to compute the Verma module of that
via universal enveloping definition.
This results in that vector space which contains the Heisenberg
algebra as $F^0$.

In fact, in section 15 of lecture notes on vertex algebras
we see that the Fock module defined in Verma-module-style as above
is in fact $\Lambda^{\infty/2}$, which is defined as
$$
\Lambda^{\infty/2}:=
\left\{\substack{\text{span of symbols of the form} \\ 
e_{i_0}\wedge e_{i_1}\wedge e_{i_1} \wedge e_{i_3}\wedge\ldots\\
\text{where }\exists N\text{ s.t. }\forall n \geq  N,
i_{n+1}=i_n+1}\right\}\Big/
\substack{\text{relation that} \\ e_i\wedge e_j=- e_j \wedge e_i\\
\text{``wherever it occurs''}}.
$$
It looks like ``in the infinite-dimensional version
we lost $q$ in the relation we quotient by''.
Indeed, when $q=0$ the Clifford algebra is just the exterior power.
\end{enumerate}

\noindent
The question here is:
why would it make sense that 
the homology of any moduli space
looked like this thing?
The first step is to look for antisymmetry,
some super-commutativity relation.

In fact, I think that \cite{Nakajima-lectures}
main theorem already shows that
infinite-dimensional Clifford is represented
in the Hilbert scheme of points.
This just corresponds to the case when
the product of the degrees of the two things to be commuted
is not even. In other words, when ``the cohomology of the variety is {\it not} 
concentrated in even degrees'',
we must have a Clifford (super) relation.

\subsection{Character formula}
\label{subsection-character}

\noindent
Here's a possible research thread left over from
my end-of-course project on 
Nakajima's representation of the Heisenberg algebra.

Recall that Nakajima constructs a representation
of the infinite-dimensional Heisenberg algebra 
on $\mathcal{H}=\bigoplus_{n\geq 0}H_*(X^{[n]})$,
where $X$ is (projective) complex surface.

Independent of the representation,
Göttsche proved the following formula
\begin{equation}
\label{equation-gottsche}
\sum_{n=0}^\infty q^nP_t(X^{[n]})
=\prod_{m=1}^\infty \frac{
(1+t^{2m-1}1^m)^{b_1(X)}
(1+t^{2m+1}q^m)^{b_3(X)}}
{(1-t^{2m-2}q^m)^{b_0(X)}
(1-t^{2m}q^m)^{b_2(X)}
(1-t^{2m+2}q^m)^{b_4(X)}},
\end{equation}

\noindent
where the {\it Poincaré polynomials} $P_t(Z)$ are defined by
$P_t(Z)=\sum_{m\geq 0}b_m(Z)t^m$, $b_m(Z)=\text{rank}H_m(Z)$ for
any variety $Z$.

Nakajima uses Göttsche's formula to prove
that his representation of the Heisenberg algebra
is in fact a highest weight representation.
Jethro's concern is that this actually looks
more ``obvious'' (once proved Nakajima's theorem).

What's the nuance? 
Couldn't there be several copies of the irreducible highest weight
representation of the Heisenberg algebra in $\mathcal{H}$?

If we could prove that $\mathcal{H}$ 
is a highest weight representation by other means,
we would be proving Göttsche's
formula in a fine way!

\subsection{Moduli space of sheaves of rank 1}
\label{subsection-moduli}

\noindent
Consider the following idea of Jethro: embed 
$$
\mathcal{H}=\bigoplus_{n\geq 0}H_*(X^{[n]})
\hookrightarrow 
\bigoplus_{\alpha \in \text{lattice}}\mathcal{H}^\alpha
$$
where $\alpha\neq 0$ are ``instanton moduli spaces''.
We have to give this a better look
(see \cite[Section 9.4]{Nakajima-lectures}),
but this direct sum of moduli spaces is a lattice vertex algebra.
Of course $\mathcal{H}$ is also a vertex algebra 
--- due to Nakajima's representation.

Morally, we would pass to the vertex algebra structure
(of $\mathcal{H}$ or of the instanton vertex algebra),
do computations there and then go back to the usual
Heisenberg algebra.

\medskip\noindent

Here's a few key points from \cite[Section 9.4]{Nakajima-lectures}.

Start with the Néron-Severi lattice $\text{NS}(X)$,
where $X$ is now also simply connected
(mostly likely it's a projective complex surface).
This has a non-degenerate symmetric bilinear pairing,
the intersection form $\left<\cdot,\cdot\right>$ 
which by the Hodge index theorem has rank $(1,n)$.

This means we can create its associated vertex algebra.
Just as the Fock module of a general vertex algebra
associated to a lattice $L$ is $V_L$,
$$
V_{\text{NS}(X)})
\cong \bigoplus_{c_1 \in\text{NS}(X),ch_2}H_{\text{NS}}(\mathcal{M}(c_1,ch_2))
$$
where $H_{\text{NS}}(\mathcal{M}(c_1,ch_2))$
is a subspace of $H_*(\mathcal{M}(c_1,ch_2))$.

\medskip\noindent
Now fix $c_1 \in L$ (I think $L$ is $\text{NS}(X)$)
and $\text{ch} \in 1/2\mathbb{Z}$.
Then we have $\mathcal{M}(c_1,\text{ch}_2)$,
the moduli space of rank 1 torsion-free sheaves
with $c_1(E)=c_1$.

Can you find the charged free fermions anywhere around here?

\subsection{Charged free fermions}
\label{subsection-fermions}

\noindent
The basic fact about the fermions is that they
are the Fock module of the Lie algebra
$$
\tilde{\mathfrak{a}}=\underbrace{t^{1/2}\mathfrak{a}[t,t^{-1}]}_{\text{odd}}
\oplus \underbrace{\mathbb{C}K}_{\text{even}}
$$
where
$$
\mathfrak{a}=\mathbb{C}\varphi +\mathbb{C} \varphi^*,
$$
and
\begin{align*}
\left<\cdot,\cdot\right>: \mathfrak{a}\times \mathfrak{a}
&\longrightarrow  \mathbb{C}\\
\left<\varphi^*,\varphi\right> &=1\\
\left<\varphi,\varphi^*\right>&=1
\end{align*}

\noindent
The Heisenberg algebra $H$
is contained in the charged free fermions $F$,
$$
\underbrace{H}_{\left<\alpha(w)\right>}
=F^{(0)}
\subset 
\underbrace{F}_{\left<\varphi(w),\varphi^*(w)\right>}
=\bigoplus_{m \in \mathbb{Z}}F^{(m)}.
$$
where $\alpha= :\varphi\varphi^*:$,
the normally ordered product.

Recall that the charged free fermions have character
\begin{align*}
\text{Tr}_Fy^{\alpha_0}q^{L_0}
&=\chi_F\\
&=\prod_{p=1}^\infty
(\underbrace{1-yq^{p-1/2}}_{\varphi^*-p+1/2})
(\underbrace{1-y^{-1}q^{p-1/2}}_{\varphi-p+1/2})\\
&=\prod_{m=1}^\infty
\frac{1}{1-q^m}
\sum_{k \in \mathbb{Z}}y^kq^{k^2/2}.
\end{align*}

{\bf This makes me think}
that Jethro may want to find a representation
of the charged free fermions in the
cohomology of the moduli spaces.

\noindent
Now it gets even more dreamy:
``There is some physicist who, for physics reasons…''
\begin{align*}
\text{Tr}_H q^{L_0}t^{w_0}&\\
w&=:\alpha \alpha \alpha:\\
\text{Tr}_Fy^{\alpha_0}q^{L_0}t^{w_0}
&=\prod_{p=1}^\infty(1-yq^{p-1/2}t^{(p-1/2)^2})(1- \ldots),\\
[\alpha_0,\varphi_p]&=\varphi_p\\
[L_0,\varphi_p]&=p\varphi p \\
[w_0,\varphi_p]&=p^2\varphi p.
\end{align*}

\noindent
That's my best understanding for the time being.
Of course there are several hints
of what these symbols mean on \texttt{vertex-algebras.pdf} at
\url{https://github.com/danimalabares/vertex-algebras}. 


\subsection{Quiver varieties}
\label{subsection-quiver-varieties}

\noindent
From \cite{Nakajima-ALE}: ``The quiver varieties
will be described as hyper-Kähler quotients 
(see Hitchin et al. [HKRLR] of representation
spaces of a quiver on a graph,
which are finite-dimensional quaternion vector spaces.''

Here is \cite[Theorem 10.14]{Nakajima-ALE}:

\begin{theorem}
\label{theorem-rep-ale}
The operators $E_k,F_k,H_k$ on $L(\mathbf{w})$ 
give the irreducible highest-weight integrable representation
of the Kac-Moody algebra associated with the
generalized Cartan matrix $\mathbf{C}$ with the highest
weight $\mathbf{w}$.
Each summand of the decomposition 
$L(\mathbf{w})=\bigoplus_{\mathbf{v}}L(\mathbf{v},\mathbf{w})$
is a weight space with the weight $\mathbf{w}-\mathbf{C}\mathbf{v}$.
\end{theorem}


\subsection{2025 talk}
\label{subsection-nak2025}

\noindent
Here are some notes of what I picked up from
Nakajima's talk on 19/11/2025.

\medskip\noindent
Quiver varieties and representations of Yangi

They are a kind o moduli space of sheaves on $\mathbb{P}^2$.
Middle-degree cohomology is an irreducible finite-dimensional
representation of $\mathfrak{g}$ with highest weight $\omega$.
See [N '94]. An actionon the base induces an action on the moduli spaces…
This gives equivariant cohomology, which is a deformation of 
ordinary deformation.
Quantum loop algebra $U_\mathfrak{g}(L\mathfrak{g})$
is a defomration of usual loop algebra.
There's also Yangian, which a deformation quantization
of $C\mathfrak{g}=\mathfrak{g}[z]$. I think equivariant cohomology
of the quiver variety are representations of Yangian
These are more interesting than representations of the usual one:
they are no longer semisimple, and they are monoidal
(have tensor product of representations is rep)-- instead
of commutative! as the usual case.

{\bf Maulik-Okunkov stable envelope.}
A tool for understanding representation theory of quiver variety.
Decompose $\omega=\omega_1+\omega_2$.
This gives splitting of moduli space and equivariant cohomologies.A

equivariant cohomology of a point is a polynomial space
of as many variables a the dimension of the torus acting!

Choosing $v=0=2$ it becomes moduli of vector bundles!

\subsection{bandoleiros}
\label{subsection-}

work by Jimbo-Kasiwara-Miva, 1981

A simple case would be $r=1$, $B_{1}=?$.

Start with a fd vector space $V_n$, 
which is a quotient of polynomial ring of dimension $n$,
and consider $\Lambda^r V_n=$

\begin{theorem}[-,Salehyan, 2018, BSBM]
\label{theorem-salehyan}
Some operators $\sigma$ can be extended to
$\Lambda^{\infty}V\cong B=\mathbb{Q}[x_1,x_2,…]$
\end{theorem}

\begin{proof}
 $B$ has basis
parametrized by partitions, just like $\Lambda^{\infty}$!
Such a basis is a basis of Schurs polynomials at least in 
the finite-dimensional case where $B=B_r=\mathbb{Q}[x_1,…,x_r]$.
\end{proof}

\noindent
The Boson-Fermion correspondence is
\begin{align*}
: B_r &\longrightarrow \Lambda^rV \\
S_{\underline{\lambda}(\underline{x}}) &\longmapsto 
S_{\underline{\lambda}}(\underline{x})x^r(0)=x^r(\underline{\lambda}).
\end{align*}
\noindent
which induces
$$
B_{r,n}\xrightarrow{\sim}\Lambda^rV_n.
$$
In fact, $B_{r,n}\cong H^*(\text{Gr}(r,n),\mathbb{Q})$.
This makes $B_{r,n}$ into a representation of
$\mathfrak{gl}_n(\mathbb{Q})$. This is DJKM in finite dimensional case.

An observation by Renato et al. is that
$B_{1,r}\cong H^*(\mathbb{P}^?,\mathbb{Q})$

Fermions are endomorphism of exterior power.
It's generated by
$$
(x^{i_1}\wedge… \wedge x^{i_k})\otimes(\partial^{j_1}… \partial^{j_k})
$$
You can also think of it as
$V \oplus V^*$, where the duals act by
contraction and the wedges by…

\subsection{Torus action on Hilbert scheme}
\label{subsection-torus}

\noindent
From Michele's last talk

\begin{align*}
\mathbb{C}^* \times \mathbb{A}^n &\longrightarrow \mathbb{A}^n \\
(\lambda,(x_1,…,x_n) &\longmapsto (\lambda^{\mu_1}x_1,…,\lambda^{\mu_n}x_n),
\qquad \mu_i>0.
\end{align*}

\noindent
This gives an action on the polynomial ring
$\mathbb{C}[x_1,…,x_n]$, which in turn gives
an action on ideals. But the Hilbert scheme parametrizes ideals!
This gives an action $\mathbb{C}^* \mathbb{y} \text{Hilb}\mathbb{A}^n$.

\begin{definition}
\label{definition-bb-decomposition}
The Birula (BB) decomposition of $\text{Hilb}\mathbb{A}^n$ is
$$
\text{Hilb}^+\mathbb{A}^n
=\{[Z_i]_{i=1}^r \in \text{Hilb}\mathbb{A}^n
:\exists  \lim_{\lambda \to \infty}\lambda[Z_i]_{i=0}^r \}.
$$
\end{definition}

\noindent
There is a bijection (not an isomorphism! RHS is connected
and LHS isn't)
$$
\text{Hilb}^+(\mathbb{A}^n) \to (\text{Hilb}\mathbb{A}^n)_0
$$
where the RHS is the punctual hilbert scheme -
the subschemes supported at the origin of $\mathbb{A}^n$.

This gives that the fixed point set
of the torus action is ``equal''
in the BB decomposition and in the punctual hilbert scheme.






\section{Fanosearch}
\label{section-fanosearch}

\noindent
Perhaps the main objective of Fanosearch
project is to classify Fano 4-folds using
mirror symmetry (see \cite{mirror-fano}).

\medskip\noindent
A premilinary work of Galkin is \cite{small-toric}.
The main idea here taking a toric fano threefold $X$
(these are classified so we know how $X$ can be),
smoothing it (it has been shown they can be smoothed),
get the invariants of the smoothing $Y$
(which also falls in the classification),
and then show that if any smooth Fano threefold $Y$
can be degenerated to a toric Fano,
it must have been constructed as above.
We thus know all the possible $Y$.

\medskip\noindent
The problem for fourfolds is of course more
complicated since we don't have a classification.
Thre appears to be too much data
which might be understood using IA.
Possible first steps:

\begin{itemize}
\item Understand what is the data.
\item Understand how the data was constructed
(i.e. understand the theory, possibly not a good idea).
\item Try to do some hands-on computations to understand
what we are supposed to do. (Need guidance for this.)
\end{itemize}

\noindent
Also recall:

\begin{quotation}
My post from February 14, 2012 on our blog is
basically an actionable instruction for a PhD
student who can program, etc.
\end{quotation}

\noindent




\bibliography{my}
\bibliographystyle{amsalpha}

\end{document}
